\subsection*{一步一步建立直观}

\paragraph{1. 先识别什么是动作的输入}

把左图的绿色箭头当成写在弹性网格上的一段位移。输入箭头改变时，矩阵保持固定，右图显示同一规则给出的输出。若是恒等矩阵，两边坐标完全一致。这一步把改变输入与改变变换区分开来，避免把两种操作混在一起。



\paragraph{2. 两根基箭头如何决定整张网格}

剪切时第一根基箭头不动，第二根向右倾斜。任意输入 $(x,y)$ 都由 $x$ 份第一根箭头和 $y$ 份第二根箭头组成，所以输出为 $(x+y,y)$。网格不是被逐点任意拉动：所有点共享同一个线性组合规则，平行线因而仍平行或塌缩。



\paragraph{3. 压扁揭示了线性规则的边界}

把两列变成平行方向后，单位方形的面积归零，不同输入可以落到同一个输出。这不破坏线性，却破坏可逆性。纸笔证明应分别检查相加与数乘保持；不能因为动画看起来平滑，就断言动作线性，更不能进一步断言它可逆。



\subsection*{定义、命题与证明}

\paragraph{定义：线性映射}

映射 $T$ 线性，是指它保持向量加法与标量乘法：$T(u+v)=T(u)+T(v)$，$T(cu)=cT(u)$。平移一般不满足这个定义。

\paragraph{命题：基向量的像决定整个映射}

设 $e_1,e_2$ 是平面的标准基。若 $T$ 线性，则对任意坐标 $x,y$，有 $T(xe_1+ye_2)=xT(e_1)+yT(e_2)$。

\begin{proof}

先用加法保持性把两项分开，再分别把标量移出 $T$，就得到结论。把 $T(e_1)$ 与 $T(e_2)$ 写成矩阵的两列，右侧正是矩阵乘向量的列组合。

因此我们不必给平面上每个点单独记录去向。不过反过来只观察两根箭头，尚不能证明未知映射线性；线性是推论所需的前提。

\end{proof}
